\par\SourceLeaf{152}
Es bleibt mir nun noch, einige Punkte, die zum Teil schon oben berührt wurden, etwas näher zu betrachten.

Man hat oben gesehen, dass P. und R. über ein Grundprinzip ganz entgegengesetzte Ansichten hatten (S. [133] und [137]. Wenn Letzterer S. [138, oben] von einem Steigen und \textit{Sinken} des \textit{Fonds} spricht, so vermute ich, dass er eigentlich nur den \textit{Ertrag} des Fonds meinte. Denn soweit die Einrichtung eine Zuwendung darstellt, scheint es mir die strenge Pflicht ihrer Verwaltung zu sein, dafür zu sorgen, dass die Substanz, aus der die Zuwendung fließt, unversehrt erhalten bleibt. Schon wegen der Verluste, die bei der Ausleihe von Kapital von Zeit zu Zeit trotz aller Vorsicht nicht zu vermeiden sind, lässt sich dies nur dann mit Sicherheit erreichen, wenn man sich neben der Erhaltung auch eine gewisse allmähliche Vermehrung zum Ziel setzt. Dabei möge man immer lieber etwas zu viel als zu wenig tun. Auf diese Weise wird die \textit{Gesamtheit} der späteren Empfänger gegenüber der Gesamtheit der früheren nicht zu kurz kommen, sondern eher besser dastehen. Die gegenwärtigen Empfänger können ihnen dies umso eher gönnen, da sie selbst die Früchte einer ähnlichen Zurückhaltung \textit{ihrer} Vorgänger genießen. Außerdem verlangt die Gerechtigkeit, dass man sich bemüht, die unvermeidlichen und bei einer kleinen Gesellschaft verhältnismäßig sehr großen Schwankungen der Empfängerzahl durch geeignete Maßnahmen so weit wie möglich auszugleichen. Dagegen scheint mir P.s Forderung, dass niemals ein künftiger \textit{einzelner} Empfänger weniger erhalten solle, als irgendein früherer erhalten hat, bei einer Gesellschaft, die, wie \textsc{Sch.} sehr richtig bemerkt hat, immer etwas Aktienmäßiges behalten wird, rechtlich nicht begründet. Jedenfalls aber, und das ist der entscheidende Punkt, kann eine solche Forderung, wenn sie als unbedingt gelten soll, überhaupt nicht ohne die offenkundigste Ungerechtigkeit gegenüber den gegenwärtigen Empfängern erfüllt werden. Es liegt auf der Hand: Je größer die Sicherheit sein soll, dass jener Fall niemals eintreten müsse, desto weniger darf man den jetzigen Empfängern geben. Man müsste die äußersten Fälle für die mögliche Empfängerzahl berücksichtigen; davon waren P.s Ansätze, wie oben gezeigt wurde, weit entfernt. Noch viel deutlicher zeigt sich dies an den weiter unten S. [150] aufgestellten Überschlägen auf Grundlage einer vergrößerten Teilnehmerzahl. Schon 1794 hätte man diese Möglichkeit zu den künftig nicht bloß möglichen, sondern sogar wahrscheinlichen Fällen rechnen können. Denn damals betrug die Gesamtzahl aller Universitätsprofessoren 45\footnote{Einer von ihnen, K., starb noch vor Abschluss der Beratungen am 1794 August 21. Jetzt (im Herbst 1845) sind es 57. Die Zahl der verheirateten Mitglieder der Witwenkasse ist aber in einem viel stärkeren Verhältnis gestiegen, von 26 im Jahr 1794 auf 42 im Jahr 1845.}, fünfzig Jahre früher nur 25. Welche Grenzen die ständig wachsenden Bedürfnisse der Zeit finden werden, lässt sich im Voraus unmöglich festsetzen.

Der zweite Punkt betrifft die Auslegung der Progressionsregelung im Sinne von \textsc{Sch.}. Wenn \textsc{Sch.} sagt [S. 138, Z. 19], P. habe das Maß der Erhöhung und ebenso der Verminderung durch eine unveränderliche Regel bestimmt, so setzt dies zwar notwendig die erste Auslegung S. [129] voraus; dennoch räumt \textsc{Sch.} damit zu viel ein. Eine unveränderliche Regel für die Verminderung lag darin nur insofern, als bestimmt wurde, \textit{wann} eine Verminderung eintreten müsse (nämlich nach jener Auslegung unmittelbar nach Überschreiten der Zahl von 15 Pensionen), aber noch nicht für die \textit{Größe} der Verminderung. Diese Unvollständigkeit hätte, wie mir scheint, nach der damaligen Sachlage am angemessensten durch eine Bestimmung in folgender Fassung ergänzt werden können:
\begin{quote}
Jedes Mal, wenn das Kapitalvermögen um 5000 Rthl. gestiegen ist, tritt eine Erhöhung der Pensionen ein, die für jede einzelne Pension 10 Rthl. beträgt, wenn und solange nicht mehr als 15 Pen%
\SourceLeaf{153}sionen bestehen, sonst aber insgesamt 150 Rthl., die gleichmäßig auf die einzelnen Pensionen zu verteilen sind. Dasselbe gilt für jede folgende Erhöhung.
\end{quote}
\noindent Auf diese Weise hätte man nicht einmal den \textit{Beginn} einer Erhöhung davon abhängig machen müssen, dass die Zahl der Pensionen nicht über 15 hinausgeht. In dieser Hinsicht hätte sich diese Art der Bestimmung sogar als vorteilhafter für die Witwen erwiesen. Denn bei einer andauernden kleinen Überzahl würde der Verlust, der aus der Verpflichtung entsteht, die bisherige Erhöhungssumme mit mehreren teilen zu müssen, bald durch eine neue Erhöhung ausgeglichen oder mehr als ausgeglichen, während in einem solchen Fall eine neue Erhöhung nach P.s Regelung und in der zweiten Auslegung überhaupt nicht zulässig ist.

Im Übrigen möchte ich glauben, dass die obige geänderte Fassung, wenn damals jemand daran gedacht hätte, sie vorzuschlagen, auch P. hätte zufriedenstellen müssen. Denn entweder war seine Voraussetzung, 15 sei die höchste mögliche Pensionenzahl, richtig oder sie war unrichtig: Im ersten Fall blieb die Änderung ganz wirkungslos und war somit gleichgültig, im zweiten war sie notwendig.

Die Folge einer solchen Regelung wäre gewesen, dass man sich daran gewöhnt haben würde, die Pension als aus zwei Teilen zusammengesetzt zu betrachten: einem unveränderlichen Teil und einer Zulage, die von Zeit zu Zeit mit dem Kapitalvermögen wachsen, möglicherweise aber dabei auch wieder etwas zurückgehen könnte. Letzteres geschähe dann jedoch nach einer wirklich unveränderlichen einfachen Regel, anhand derer jede Witwe die Kontrolle leicht selbst führen könnte. Auf den Unterschied zwischen einem solchen regelgemäßen Zurückgehen und demjenigen, das nach der Quittungsklausel infolge des Unvermögens der Kasse eintritt, habe ich schon oben S. [131] aufmerksam gemacht.

